% 3-3-2-live-data.tex
%  Budget: 2pp.  The OOD / tool-only dataset.
%  Source map: no external literature claims in this subsection; it documents
%   our own data path. Attribution line is contractual (CC BY 4.0), not a cite.
%  REQUIRED verbatim attribution (configs/data.yaml):
%   "Data: Energy-Charts (Fraunhofer ISE) / Bundesnetzagentur SMARD.de, CC BY 4.0"
%  Numbers: 173 complete days of 2026, 4343 hourly rows, zone DE-LU.
%  PRE-DISCLOSE the definition mismatch here (HANDOFF 6.8 amendment) -- it is
%  the mechanism behind chapter 5's negative result and must not surprise
%  the reader there.

\subsection{داده زنده \lr{Energy-Charts}}
\label{sec:3-3-2}

\subsubsection*{دامنه مجاز استفاده}

مجموعه‌ی دوم این پژوهش از سامانه‌ی \lr{Energy-Charts} مؤسسه‌ی \lr{Fraunhofer ISE}
دریافت شده است. کاربرد آن به‌صراحت محدود است و این محدودیت باید پیش از هر توصیفی
ثبت شود: این داده تنها در آزمون برون‌توزیعیِ \cref{sec:5-2} و در ابزار عملیاتیِ \cref{sec:5-3} به
کار می‌رود، و هیچ عددی از آن وارد جدول‌های نتایج فصل چهارم نمی‌شود.

دلیل این تفکیک، حفظ هم‌مقیاسیِ مقایسه است. تمام ادعاهای دقتِ فصل چهارم در برابر
اعداد منتشرشده‌ی مرجع سنجیده می‌شوند و این سنجش تنها بر همان مجموعه‌داده، همان دوره و
همان تعریف متغیرها معنا دارد. آمیختن دو منبع، مقایسه را در سکوت بی‌اعتبار می‌کرد.

\subsubsection*{مشخصات}

داده از ناحیه‌ی پیشنهاددهیِ \lr{DE-LU} گرفته شده و ۱۷۳ روز کاملِ سال ۲۰۲۶، معادل
\lr{4{,}343} سطر ساعتی و بدون جاافتادگی، را پوشش می‌دهد. واسط برنامه‌نویسیِ این سامانه
بدون کلید و بدون ثبت‌نام در دسترس است، که آن را با قاعده‌ی بازتولیدپذیریِ پژوهش سازگار
می‌کند. داده‌های اخیر با تفکیک پانزده‌دقیقه‌ای عرضه می‌شوند و در بارگذار پروژه به
تفکیک ساعتی تجمیع می‌گردند تا با شِمای مشترک دو مجموعه هم‌شکل شوند.

یک ناپیوستگیِ تعریفی در سطح ناحیه نیز شایان ذکر است: ناحیه‌ی پیشنهاددهی پیش از یکم
اکتبر ۲۰۱۸ \lr{DE-AT-LU} بوده و پس از آن به \lr{DE-LU} تفکیک شده است. دوره‌ی بنچمارک
پیش از این تفکیک و دوره‌ی زنده پس از آن قرار می‌گیرد.

انتشار این داده تحت پروانه‌ی \lr{CC BY 4.0} است و اسناد آن، ذکر منبع را الزامی
می‌کند. عبارت ارجاع، به شکلی که پروانه اقتضا می‌کند، چنین است:

% The licence string must appear verbatim (configs/data.yaml). At full size in
% a quote block it wrapped after "CC BY" and orphaned "4.0" on its own line,
% splitting the licence name -- bad for a compliance credit. \small plus tied
% spaces in the licence name keeps it on one unbroken line.
\begin{center}
\small\lr{Data: Energy-Charts (Fraunhofer~ISE) / Bundesnetzagentur~SMARD.de,
CC~BY~4.0}
\end{center}

\subsubsection*{یک ناهم‌خوانی تعریفی که باید همین‌جا ثبت شود}

میان این مجموعه و مجموعه‌ی بنچمارک دو ناهم‌خوانی وجود دارد که خواننده باید پیش از
رسیدن به نتایج فصل پنجم از آن‌ها آگاه باشد. ثبت‌کردن آن‌ها در همین‌جا \lr{—} یعنی
پیش از آنکه نتیجه‌ای گزارش شود \lr{—} تصمیمی آگاهانه است: این ناهم‌خوانی سازوکارِ
اصلیِ یافته‌ی منفیِ \cref{sec:5-2} است، و اگر تنها در کنار همان نتیجه معرفی شود، شبیه توجیهِ
پس از وقوع به نظر می‌رسد.

ناهم‌خوانیِ نخست جغرافیایی است. متغیر برون‌زای بار در دوره‌ی آموزش پیش‌بینی بارِ ناحیه‌ی
کنترلیِ \lr{Amprion} است، حال آنکه در دوره‌ی سرویس‌دهی پیش‌بینی بارِ ملیِ آلمان
دریافت می‌شود. سهم \lr{Amprion} از بار آلمان حدود چهل درصد است، پس دو سری زیر یک نام
در دو مقیاس متفاوت قرار می‌گیرند. اندازه‌ی دقیق این اختلاف و پیامد آن برای هر مدل در
\cref{sec:5-2} سنجیده و گزارش می‌شود.

ناهم‌خوانیِ دوم در دامنه‌ی متغیرهاست. سطح قیمت در سال ۲۰۲۶ به‌مراتب بالاتر از دوره‌ی
بنچمارک است و بازه‌ی نوسان آن نیز گسترده‌تر؛ یعنی مدل‌ها در دوره‌ی سرویس‌دهی با مقادیری
روبه‌رو می‌شوند که در دوره‌ی آموزش نظیری نداشته‌اند.

هیچ‌یک از این دو، آزمون برون‌توزیعی را بی‌اعتبار نمی‌کند؛ اما هر دو معنای آن را
تغییر می‌دهند. آن آزمون، ترکیبی از تغییر واقعیِ بازار و ناهم‌خوانیِ مسیر داده‌ی خودِ
ما را می‌سنجد، و \cref{sec:5-2} این دو را تا آنجا که سنجش اجازه می‌دهد از هم تفکیک می‌کند.

\subsubsection*{بازتولیدپذیری و اسناد}

همان قاعده‌ی تصویرِ ثابت که در \cref{sec:3-3-1} برای مجموعه‌ی بنچمارک شرح داده شد، بر این
مجموعه نیز حاکم است: پنجره‌ی ۱۷۳روزه یک‌بار دریافت، با اثر انگشت \lr{sha256} ثبت و در
مخزن نگهداری شده است، و ارزیابی‌های \cref{sec:5-2} تنها از همان تصویر خوانده می‌شوند. بدون
این قید، آزمون برون‌توزیعی با هر بار اجرا عدد دیگری می‌داد و هیچ‌یک از آن اعداد
قابل استناد نبود.
